% Program W — External validation of the positive-closure proof architecture
% Normative spec: specs/IN-PROCESS/EPIC_016_WV2A_EXTERNAL_VALIDATION_BENCHMARK_PAPER_SPEC.md
\documentclass[11pt]{article}
\input{../suite_preamble}
\usepackage{longtable}
\usepackage{listings}
\lstdefinelanguage{lean}{
  keywords={def, theorem, lemma, structure, class, instance, where, by, exact, rfl, simp, intro, have, let, refine, rcases, obtain, cases, fun},
  keywordstyle=\bfseries,
  ndkeywords={Prop, Type, Sort, Nat},
  ndkeywordstyle=\bfseries,
  sensitive=true,
  comment=[l]{--},
  morecomment=[s]{/-}{-/},
  commentstyle=\itshape,
  stringstyle=\ttfamily,
  basicstyle=\ttfamily\small,
  breaklines=true,
  showstringspaces=false
}
\input{../suite_macros}

\newcommand{\IC}{\ensuremath{\mathsf{IC}}}

\title{Program W: External Validation of the Positive-Closure Proof Architecture\\[0.35em]
\large Quotients, Fibers, and Canonical Sections in Lean}
\author{Nova Spivack}
\date{2026}

\begin{document}

\maketitle

\begin{abstract}
The Infinity Compression (\IC{}) formalization develops a layered \emph{positive-closure} proof architecture: collapse to a bare certificate, a typed forgetful map, fiber structure over that certificate, and---where appropriate---canonical witnesses or sections. A legitimate methodological concern is that this architecture might be an artifact tailored to the flagship NEMS-style spine rather than a reusable proof pattern. \textbf{Program W} is the first machine-checked \emph{external} validation tranche: we formalize an ordinary quotient/setoid benchmark in Lean, package it in the same collapse/forgetful/fiber/section language, and compare a flat baseline (\texttt{Quotient.exists_rep}) to an architecture-guided theorem bundle (nonempty fibers, surjectivity of the projection, existence of a global right inverse, Mathlib's \texttt{Quotient.out} as a canonical section). We record an explicit validation verdict (\textbf{STRONG}) in the EPIC\_015\_WV1 sense: not new classical mathematics, but a strengthened, organized interface that exports cleanly from the \IC{} story. This paper states the honest scope---external validity and transferability, not a universal claim over all future closure problems---and positions Program W within the larger NEMS / APS / IC / Strata stack.
\end{abstract}

\tableofcontents
\newpage

\section{Introduction}
What is being validated; why internal success of the flagship formalization is insufficient; why an external benchmark matters for methodology.

\section{The positive-closure architecture}
Collapse; forgetful projection; fiber; section or canonical witness. Brief alignment with the flagship IC paper and EPIC\_009 fiber language.

\section{Benchmark choice and baseline}
Setoid quotients; classical representative existence; baseline theorem profile (\texttt{exists_rep}).

\section{Lean formalization}
Path: \texttt{InfinityCompression/Validation/QuotientFiberBenchmark.lean}. Summary of \texttt{BareQuotient}, \texttt{EnrichedCarrier}, \texttt{forgetToQuotient}, \texttt{QuotientFiber}; theorem bundle including \texttt{forgetToQuotient\_hasRightInverse}, \texttt{quotientOut\_rightInverse\_forgetToQuotient}, \texttt{canonicalQuotientFiberWitness}.

\section{What the architectural packaging adds}
Right inverse and \texttt{HasRightInverse}; canonical witness; section language; comparison to flat baseline.

\section{Validation verdict and epistemology}
Why \textbf{STRONG} in the EPIC sense; why this is not a claim of new pure quotient theory; methodological significance.

\section{Place in the larger program}
NEMS (outer burden); APS (local-to-global exactness); IC flagship (internal positive closure); Program W (external validation); Strata (abstract synthesis). Program W answers the concern that the positive architecture is only endogenous to \IC{}.

\section{Limitations and future work}
Single external benchmark to date; deferred multi-benchmark survey paper (EPIC\_016\_WV2B); optional harder field-facing targets.

\appendix
\section{Theorem index (Lean names)}
% Mirror \texttt{MANIFEST.md} and EPIC\_015\_WV1 inventory; map statements to \texttt{QuotientFiberBenchmark.lean}.

\begin{thebibliography}{9}
\bibitem{EPIC015} EPIC\_015\_WV1 specification and Program W notes (repository).
\end{thebibliography}

\end{document}
